\documentclass[11pt,a4paper]{article}
\usepackage[margin=2.2cm]{geometry}
\usepackage[T1]{fontenc}
\usepackage{lmodern}
\usepackage{xcolor}
\usepackage{graphicx}
\usepackage{booktabs}
\usepackage{tabularx}
\usepackage[hidelinks,colorlinks=true,linkcolor=febblue,urlcolor=febblue]{hyperref}
\definecolor{febblue}{HTML}{1C4687}
\definecolor{febgold}{HTML}{FDB515}
\newcommand{\file}[1]{\texttt{#1}}
\setlength{\parskip}{4pt}

\title{\color{febblue}\LARGE Autonomous Racing Hardware Proposal\\[4pt]\large Formula Electric at Berkeley, Autonomous Team}
\author{Prepared for the FEB leadership and budget committee}
\date{September 2026}

\begin{document}
\maketitle

\section*{Summary}

Over the last month the autonomous team built and shipped a complete sim-racing platform:
a native simulator of the RoboRacer 1:10 race car, a containerised ROS~2 devkit, a starter
driving stack, a competition harness and a public leaderboard, with a two-stage season
(qualification, then a head-to-head final) and written manuals for competitors and organisers.
Members can be driving a car in software within an hour of joining. What they cannot do yet is
put that software on a real car. This proposal asks for \textbf{one RoboRacer-standard 1:10
vehicle and a track kit now (\$3{,}500 including contingency) and a second car plus a camera
next semester (\$2{,}900)}, so that
the same code members write for the simulator runs on hardware, and so FEB can enter the
RoboRacer physical competitions.

\section{What we have built}

\begin{tabularx}{\linewidth}{lX}
\toprule
Piece & What it does \\
\midrule
FEB Simulator & Native Mac, Windows and Linux app; a fork of the AutoDRIVE simulator carrying the RoboRacer digital twin unchanged. Tracks load from folders; head-to-head with up to four cars; a house racer to practise overtakes against. \\
Devkit & A Docker image (ROS~2 Humble) that turns the simulator's sensor stream into ROS topics and hosts the member's code. Foxglove for visualisation on any laptop. \\
Starter kit & A commented reactive driver that laps any track, with a tuning file and a system-identification tool; the learning path goes from tuning it to mapping, racelines and model-predictive control. \\
Competition harness & Runs every team's submitted image on a secret track under league rules, records bags, audits banned topics, publishes results. \\
Leaderboard site & Practice boards per track, qualification status per team, competition standings and brackets. \\
Manuals & A competitor manual (install to leaderboard entry) and an organiser manual (tracks, events, race day, releases). \\
\bottomrule
\end{tabularx}

\textbf{Onboarding.} A new member clones one repository, runs one setup command, and drives
manually. Copying the starter kit and running one more command puts an autonomous car on
track. Their first leaderboard entry follows the same afternoon. No ROS installation, no
hardware, no organiser time.

\textbf{The season.} Members qualify by completing ten laps with at most one collision on
every practice track. Qualified teams race a time-attack on a secret track; the top seeds meet
in a two-car knockout final. All of this exists and has been exercised end to end.

\section{Why hardware, and why now}

\begin{itemize}
  \item \textbf{Sim-to-real is the point of the exercise.} The simulator's car is the digital twin of the RoboRacer reference vehicle. A stack tuned in it is meant to transfer to the real car; without one, that claim is never tested and members never meet latency, tyre variation, or a lidar that sees the room.
  \item \textbf{Competitions.} RoboRacer (formerly F1TENTH) runs physical races at ICRA, IROS and IFAC each year, on the same rules our league uses. Our simulator entries already match the league's format; a car makes the physical events reachable.
  \item \textbf{Retention.} Software-only teams lose members after the first semester. A car on a table is what keeps a room full.
\end{itemize}

\section{Platform choice}

\subsection{Traxxas Slash 4x4 (1:10), the RoboRacer standard}

The RoboRacer build is defined around the Traxxas Slash 4x4 Premium chassis. Everything
downstream assumes it: the published bill of materials and build videos, the laser-cut
mounting plates, the VESC motor-controller settings, the tyre and wheelbase numbers in the
simulator, and the rules of every physical competition (chassis codes TRA74054, TRA6804R,
TRA68086). Choosing it means the team inherits a community's solved problems instead of
solving them alone.

\subsection{Why not the team's Tamiya 1:10 chassis}

A motorless Tamiya touring-car kit is available. It would save the chassis cost (about \$350)
and cost a semester: different wheelbase and track width from the simulator model, no
brushless motor or controller pairing to copy, no sensor plates, and no competition
eligibility. That is mechanical engineering, not autonomy, and it is work no one else has
documented. Recommendation: keep it as a bench chassis for practising sensor mounting and
wiring, and race on the Slash.

\subsection{Why not a larger scale}

A 1:5 or larger car needs an outdoor-sized track, carries several times the energy, cannot be
run safely in a lab, and has no autonomous-racing community around it. Traxxas makes no
road-going car at that scale in any case. Scale 1:10 is where the league, the parts, the
tracks (33~cm air-duct walls in a lab) and our simulator all are.

\section{Bill of materials, one car}

Prices are September 2026 list prices in USD, rounded; verify at order time.

\begin{tabularx}{\linewidth}{Xrl}
\toprule
Item & Est. & Note \\
\midrule
Traxxas Slash 4x4 Premium chassis (TRA74054) & 380 & the RoboRacer reference chassis \\
Velineon 3500 brushless motor & 70 & if not included with the chassis \\
VESC 6 MkVI motor controller & 260 & the standard controller; has an IMU \\
NVIDIA Jetson Orin Nano developer kit & 250 & any suitable computer is allowed by the rules \\
2D lidar: Hokuyo UST-10LX & 1{,}500 & the reference sensor; 40~Hz, 270\textdegree, 10~m \\
\quad \emph{alternative:} Slamtec RPLidar S2 & (400) & saves \$1{,}100; slower and noisier, still legal \\
Power: two 3S 5000~mAh LiPo, charger, LiPo bag & 200 & \\
Power distribution board (RoboRacer powerboard) & 100 & \\
Mounting plates, standoffs, cables, hub, fasteners & 120 & laser-cut from the published drawings \\
Spares: tyres, pinion, bumper & 80 & \\
\midrule
\textbf{Total with Hokuyo} & \textbf{2{,}960} & \\
\textbf{Total with RPLidar S2} & \textbf{1{,}860} & \\
\bottomrule
\end{tabularx}

\medskip
The recommended purchase is \textbf{the Hokuyo configuration} for the first car: it is what
the league and our simulator model, and the lidar is the one part that decides how fast the
car can safely go. If the budget forces a choice, buy the RPLidar car now and upgrade the lidar
for competition season.

\section{Track kit}

\begin{tabularx}{\linewidth}{Xrl}
\toprule
Item & Est. & Note \\
\midrule
Flexible aluminium air duct, 33~cm diameter, 30~m & 150 & the league's track walls; folds flat for storage \\
Blue and yellow cones, floor tape, lap-timing gate (photo interrupter) & 90 & \\
\midrule
\textbf{Total} & \textbf{240} & \\
\bottomrule
\end{tabularx}

\section{The ask, in two phases}

\begin{tabularx}{\linewidth}{lXr}
\toprule
Phase & Contents & Total \\
\midrule
A, this semester & One car with Hokuyo lidar, track kit, 10\,\% contingency & \$3{,}500 \\
B, next semester & Second car (RPLidar, for head-to-head and as a spare), Intel RealSense D435i depth camera for the cone-perception project, spares & \$2{,}900 \\
\bottomrule
\end{tabularx}

\medskip
Phase A alone gives every qualified team a car to run on. Phase B is what makes the
head-to-head final and the perception project real.

\section{Timeline}

\begin{tabularx}{\linewidth}{lX}
\toprule
When & Milestone \\
\midrule
Order + 2 weeks & Car assembled from the published build guide; manual drive; sensors on ROS~2 \\
+ 4 weeks & Starter kit laps the lab track; simulator-to-car gap measured and written up \\
+ 8 weeks & Qualified teams run their own stacks on the car; first internal physical time-attack \\
Spring & Entry to a RoboRacer physical competition \\
\bottomrule
\end{tabularx}

\section{Risks}

\begin{itemize}
  \item \textbf{Crashes.} The Slash is a bashing truck; the lidar is the fragile part. Bumper, walls of soft duct, and a speed limit for new stacks. Spares are budgeted.
  \item \textbf{Battery safety.} LiPo bag, supervised charging, a written procedure before the first charge.
  \item \textbf{Supply.} Hokuyo lead times vary; order it first. Every other part is stocked by hobby and electronics retailers.
  \item \textbf{Ownership.} The car lives with the autonomous team lead; access is through the same qualification rule as the simulator, so it is never an untested stack's first target.
\end{itemize}

\bigskip
\noindent Simulator, devkit, manuals and results: \url{https://github.com/Pranman1/feb-racing}\\
Leaderboard: \url{https://pranman1.github.io/feb-racing/}\\
RoboRacer build guide and bill of materials: \url{https://f1tenth.readthedocs.io}\\
RoboRacer competition rules: \url{https://github.com/f1tenth/roboracer_rules}

\end{document}
